% TOP_PROBLEM: {"id": 2200016, "title": "Problems Around Polynomials — Problem 6", "category_id": 4, "category": {"id": 4, "name": "algebra", "display_name": "Algebra", "description": "Group theory, ring theory, field theory, and algebraic structures.", "slug": "algebra", "order_index": 4, "created_at": "2026-07-31T15:26:25.671Z"}, "status": "open", "problem_number": "AMR-021-0016", "set_id": 13, "statusQualification": "The missing Rolle inequalities and polynomial-like convention are restored from Section VIII."}
% FIRST_POSTED: 2026-10-01
% ENTRY_KIND: statement_only
% UnsolvedMath ID 2200016; source catalogue code AMR-021-0016
% !TeX program = lualatex
% Definition and sources only; this file is not an LLM solution attempt.
\documentclass[11pt,a4paper]{article}
\usepackage[margin=25mm]{geometry}
\usepackage{fontspec}
\setmainfont{lmroman10-regular.otf}[BoldFont=lmroman10-bold.otf,
 ItalicFont=lmroman10-italic.otf,BoldItalicFont=lmroman10-bolditalic.otf]
\usepackage{amsmath,amssymb,mathtools}
\usepackage{unicode-math}
\setmathfont{latinmodern-math.otf}
\AtBeginDocument{\let\setminus\smallsetminus}
\usepackage{xurl,hyperref}
\hypersetup{colorlinks=true,urlcolor=blue}
\setcounter{secnumdepth}{0}
\setlength{\parindent}{0pt}
\setlength{\parskip}{5pt}
\setlength{\emergencystretch}{3em}
\begin{document}
\section*{Problems Around Polynomials — Problem 6}\label{2200016}
{\small UnsolvedMath ID 2200016 · AMR-021-0016 · Algebra · AMR Open Problem Lists}
\subsection{Definitions and mathematical statement}
Let $n\geq1$ and let $I\subset\mathbb R$ be an open interval.
A smooth function $f:I\to\mathbb R$ is polynomial-like of degree $n$
if $f^{(n)}$ never vanishes on $I$. It is real-rooted if it has
$n$ distinct zeros in $I$. In that case repeated use of Rolle's
theorem shows that $f^{(i)}$ has exactly $n-i$ zeros in $I$, all
simple, for $0\leq i<n$.
Write these zeros in increasing order as
$x^{(i)}_1<\cdots<x^{(i)}_{n-i}$.

Suppose instead that the numbers in such a triangular array are
prescribed. The necessary interlacing conditions are
\[
 x^{(i)}_\ell<x^{(i+1)}_\ell<x^{(i)}_{\ell+1}
 \quad(0\leq i<n-1,\ 1\leq\ell<n-i).
\]
They imply all the longer-range inequalities
$x^{(i)}_\ell<x^{(j)}_\ell<x^{(i)}_{\ell+j-i}$ whenever the indices
are defined. Coincidences between nonadjacent derivative levels are
allowed if these strict inequalities permit them.

Characterize precisely the arrays in $I$ satisfying these conditions
that arise from some real-rooted polynomial-like function $f$ of
degree $n$. Equivalently, determine the additional compatibility
conditions among the zero locations of successive derivatives.
The realizing function need not be a polynomial, and the numerical
locations, rather than only their ordering, are prescribed.
Multiplying $f$ by a nonzero constant preserves its array, so one
may normalize the sign by requiring $f^{(n)}>0$ throughout $I$.
\subsection{Short English statement}
Which interlacing triangular arrays are exactly the zeros of a smooth function and its successive derivatives?
\subsection{Sources}
\begin{itemize}
\item[S1] ULAM.AI and the UnsolvedMath Contributors, supplied problems.json, record 2200016 (AMR-021-0016), AMR Open Problem Lists. Catalogue identity and source transcription; CC BY 4.0.
\url{https://huggingface.co/datasets/ulamai/UnsolvedMath}.
\item[S2] Shapiro - Problems Around Polynomials: The Good, The Bad and The Ugly (2015), source-order item 16 (Problem 6 as printed). Original source indexed by AMR-021-0016.
\url{https://arxiv.org/abs/1503.05295}.
\end{itemize}
\end{document}
